\documentclass[10pt,titlepage]{jsarticle}
\usepackage[utf8]{inputenc}

\usepackage[margin=23truemm]{geometry}
\usepackage{wrapfig}
\usepackage[dvipdfmx]{graphicx}
\usepackage[dvipdfmx]{color}
\usepackage{booktabs}

\usepackage{colortbl}
\usepackage{multirow}
\usepackage{tabularx}

\usepackage{enshuc}

\title{情報科学演習C \\ 課題1} 
\author{阪大太郎}
\studentid{0123456789}
\affiliation{基礎工学部xxコース}
\date{2026年x月xx日} 

\begin{document}

\maketitle 

\textcolor{blue}{青字の注意書きは、提出の際には削除すること。}

\textcolor{blue}{実行結果を載せるときは，実際にやったことが分かるように，terminalの出力の結果やそのスクリーンショットを示しつつ，説明すること}

\section{課題1-1}


\subsection*{1. \texttt{ping}コマンドについて}

\subsection*{2. ウェブブラウザの挙動}

\textcolor{blue}{スクリーンショットを図として示した上で、なぜそうなったのか考察すること。}

\subsection*{3. \texttt{nslookup}コマンドについて}


\section{課題1-2}

\subsection*{4. \texttt{/usr/sbin/arp -a}コマンドについて}

\subsection*{5. \texttt{ping}と\texttt{/usr/sbin/arp -a}の関係}

\subsection*{6. \texttt{/usr/sbin/traceroute}について}

\subsection*{7. 演習室のネットワーク構成の予想}

\subsection*{8. \texttt{/bin/netstat -r}について}


\subsection*{9. ARPの仕組み}


\section{課題1-3}

\subsection*{10. C言語の標準ライブラリ関数とシステムコールの違い}

\subsection*{11. \texttt{strace}について}


\section{発展課題}

\subsection*{1. 各種ネットワーク関連コマンドについて}

\subsection*{2. \texttt{strace}の挙動の違い} 

\section{感想，意見，疑問等}

\end{document}
